\chapter{Калибровочно-замкнутое пространство полей и суперсвязностная развилка}
\label{ch:v8-gauge-closed-field-space-superconnection-gate}

\section{Исправление типизации}

Предыдущий гейт показал, что старый 27-мерный вещественный срез поля нельзя
отождествить с 27-мерным комплексным шумовым факторпространством. Поэтому сначала
строится не гессиан шумовых каналов, а единое пространство полей.

Пусть $E_s$ и $E_t$ --- концевые пучки рангов 11 и 10 с уже существующей
представленной калибровочной алгеброй размерности 12. Полное поле переноса является
сечением
\begin{equation}
 A\in\Gamma\bigl(\operatorname{Hom}(E_s,E_t)\bigr).
 \label{eq:v8-full-transfer-field-section}
\end{equation}
Минимальное калибровочное замыкание имеет размерность 15 над $\mathbb C$, поэтому
конечный внутренний полевой слой содержит
\begin{equation}
 2\cdot15+12=42
 \label{eq:v8-gauge-closed-field-real-dimension}
\end{equation}
вещественных компонент: 30 компонент переноса и 12 коэффициентов калибровочной
связности на каждую пространственно-временную одноформу. Новых калибровочных
генераторов и концевых фермионов не добавляется.

Инфинитезимальное действие имеет стандартный вид
\begin{equation}
 \rho_{\rm tr}(X_s,X_t)A=i(X_tA-AX_s),
 \qquad
 \rho_{\rm ad}(X)Y=i[X,Y].
 \label{eq:v8-field-space-infinitesimal-action}
\end{equation}
Обе части замкнуты. Максимальные остатки равны соответственно
$6.12\cdot10^{-16}$ и $1.49\cdot10^{-15}$, а дефекты
антиэрмитовости меньше $2.1\cdot10^{-16}$.

\section{Стандартная суперсвязностная кривизна}

На $E_s\oplus E_t$ естественный кандидат равен
\begin{equation}
 \mathbb A=
 \begin{pmatrix}
  \nabla_s&A^*\\
  A&\nabla_t
 \end{pmatrix}.
 \label{eq:v8-endpoint-field-superconnection}
\end{equation}
Его кривизна схематически содержит
\begin{equation}
 \mathbb F=
 \begin{pmatrix}
  F_s+A^*A&\nabla A^*\\
  \nabla A&F_t+AA^*
 \end{pmatrix},
 \qquad
 \nabla A=dA+\nabla_tA-A\nabla_s.
 \label{eq:v8-endpoint-field-supercurvature}
\end{equation}
На двенадцати конечных калибровочных пробах ковариантность индуцированной
производной, концевых кривизн и двух Грам-блоков выполнена с максимальными
остатками
\begin{equation}
 1.64\cdot10^{-15},\qquad
 5.26\cdot10^{-16},\qquad
 6.31\cdot10^{-15}.
 \label{eq:v8-field-supercurvature-covariance-residuals}
\end{equation}
Тем самым прежняя несогласованность типов устранена на кинематическом уровне: калибровочные и
поля переноса действительно являются частями одного ассоциированного
полевого объекта. Это соответствует общему суперсвязностному и
языку произведения Каспарова \cite{v8-field-space-roepstorff,
v8-field-space-brain}.

\section{Полярная развилка полной калибровочной ковариантности}

Нужный селектор Тома VII использовал полярную коизометрию $U$ и
относительную кривизну
\begin{equation}
 \mathcal R_U(A)=AA^*U-UA^*A.
 \label{eq:v8-field-space-relative-polar-curvature}
\end{equation}
Если $U$ фиксирована как матрица выбранного вакуума, то при полном калибровочном
преобразовании $A$ ковариантность нарушается: на двенадцати пробах дефект
лежит между $2.0831$ и $21.3772$. Такая ветвь корректна только относительно
стабилизатора $U$.

Если же преобразовывать
\begin{equation}
 U\longmapsto g_tUg_s^*,
 \label{eq:v8-moving-polar-transformation}
\end{equation}
то остаток ковариантности $\mathcal R_U$ меньше $1.31\cdot10^{-14}$.
Независимое новое поле для этого не требуется на страте постоянного ранга:
полярное разложение эквивариантно,
\begin{equation}
 \operatorname{Pol}(g_tA_0g_s^*)
 =g_t\operatorname{Pol}(A_0)g_s^*.
 \label{eq:v8-polar-equivariance}
\end{equation}

Но производный полярный фон не проходит через точку запуска. Для
$\varepsilon>0$
\begin{equation}
 \operatorname{Pol}(\varepsilon A_0)=U,
 \qquad
 \operatorname{Pol}(0)=0,
 \qquad
 \|U-0\|_{\rm HS}=\sqrt{10}.
 \label{eq:v8-polar-rank-zero-jump}
\end{equation}
Полярная часть постоянна на положительном радиальном луче и скачком исчезает
в нуле. Поэтому она не является гладкой локальной функцией поля, гессиан
которой мог бы одновременно определять запуск из $A=0$.

\section{Граница родительского действия}

Стандартная кривизна \eqref{eq:v8-endpoint-field-supercurvature} даёт
калибровочную кинетику, кинетику $A$ и Грам-потенциал. Но предыдущий точный запрет
показал, что полный самосопряжённый квадрат возвращает лишние Грам-каналы и
делает неустойчивыми все исследуемые направления. Рёберный Ходж-селектор и
относительная компонента степени два пока не получены как гладкие блоки одной
полной кривизны на новом 42-мерном внутреннем пространстве полей.

\begin{theorem}[Кинематический допуск общего пространства полей]
Пятнадцатимерный комплексный модуль переноса и двенадцатимерная физическая
калибровочная алгебра образуют единое замкнутое ассоциированное пространство полей.
Индуцированная связность и стандартная суперсвязностная кривизна полностью
калибровочно ковариантны, не требуют новых калибровочных генераторов и не добавляют
фермионных состояний.
\end{theorem}

\begin{nogocriterion}[Полярный фон не является гладким полным родителем]
Нельзя считать фиксированный $U$ полной калибровочно-ковариантной структурой: он
оставляет только стабилизатор. Нельзя также без нового механизма положить
$U=\operatorname{Pol}(A)$ во всём действии: это отображение скачкообразно на
точке нулевого ранга, где должен вычисляться исходный гессиан. Полное родительское
действие требует гладкого относительного параметра порядка или полиномиальной
замены полярного селектора.
\end{nogocriterion}

\begin{protocol}
\begin{enumerate}
 \item концевые ранги: \textbf{$11$ и $10$};
 \item физическая калибровочная алгебра: \textbf{размерность $12$};
 \item полный модуль переноса: \textbf{$15$ комплексных измерений, $30$ вещественных измерений};
 \item внутренний полевой слой: \textbf{$42$ вещественных измерений};
 \item новые калибровочные генераторы: \textbf{$0$};
 \item новые фермионы концевых секторов: \textbf{$0$};
 \item индуцированная связность на $\operatorname{Hom}(E_s,E_t)$:
 \textbf{замкнута};
 \item стандартная кривизна: \textbf{калибровочно ковариантна};
 \item фиксированный полярный фон: \textbf{только стабилизаторная ветвь};
 \item движущийся полярный фон: \textbf{ковариантен при постоянном ранге};
 \item скачок при нулевом ранге: \textbf{$\sqrt{10}$};
 \item рёберно-ходжева селектор из полной кривизны: \textbf{не получен};
 \item единый BV-гессиан: \textbf{ещё не разрешён};
 \item следующий гейт: \textbf{гладкий относительный параметр порядка}.
\end{enumerate}
\end{protocol}

Машинный сертификат сохранён в
\path{s2t_v8_gauge_closed_field_space_superconnection_gate_results.json}.

\begin{thebibliography}{9}
\bibitem{v8-field-space-roepstorff}
G.~Roepstorff,
\emph{Superconnections and the Higgs Field},
arXiv:hep-th/9801040.
\bibitem{v8-field-space-brain}
S.~Brain, B.~Mesland, W.~D.~van Suijlekom,
\emph{Gauge Theory for Spectral Triples and the Unbounded Kasparov Product},
J. Noncommut. Geom. 10 (2016), 135--206; arXiv:1306.1951.
\end{thebibliography}